\documentclass[12pt]{article}
\usepackage{amsmath, amssymb, booktabs, geometry, enumitem}
\geometry{margin=0.65in}
\setlength{\parindent}{0pt}
\setlength{\parskip}{0.1em}

\title{ECON 101: Macroeconomics \\ Quiz 5}
\author{Instructor: Muhebullah Karimzada}
\date{November 26,  2025}

\begin{document}
\maketitle



\section*{Part A — Multiple Choice (5 marks)}


% Compact spacing for MCQs
\setlist[enumerate]{itemsep=3pt, topsep=3pt}

\textbf{1.} A decrease in Canadian interest rates will:
\begin{enumerate}[label=(\Alph*)]
    \item Shift AD left
    \item Shift AD right
    \item Cause a movement up along AD
    \item Shift SRAS right
\end{enumerate}
\vspace{1em}
\textbf{2.} Which of the following would shift the short-run aggregate supply (SRAS) curve to the left?
\begin{enumerate}[label=(\Alph*)]
    \item A decrease in wages
    \item A fall in oil prices
    \item An increase in corporate taxes
    \item A technological improvement
\end{enumerate}
\vspace{1em}
\textbf{3.} The long-run aggregate supply (LRAS) curve:
\begin{enumerate}[label=(\Alph*)]
    \item Slopes upward
    \item Shifts right when the price level rises
    \item Is vertical at potential GDP
    \item Shifts left when government spending increases
\end{enumerate}
\vspace{1em}
\textbf{4.} The money multiplier equals:
\begin{enumerate}[label=(\Alph*)]
    \item $1 - r_d$
    \item $\frac{1}{r_d}$
    \item $r_d$
    \item $\frac{1}{1 - r_d}$
\end{enumerate}
\newpage
\textbf{5.} If the Bank of Canada sells government securities in the open market, the most likely outcome is:
\begin{enumerate}[label=(\Alph*)]
    \item Lower interest rates and higher AD
    \item Higher interest rates and lower AD
    \item Higher inflation and higher SRAS
    \item A decrease in the desired reserve ratio
\end{enumerate}

\vspace{1em}

\section*{Part B — Numerical / Analytical Problems (15 marks)}

\textbf{Problem 1 (7 marks) — AD--AS \& Short-Run Equilibrium}

The economy is initially at long-run equilibrium:
\[
Y^* = \$2{,}500 \text{ billion}, \qquad P = 120.
\]

A negative shock reduces autonomous investment by \$50 billion.  
Assume MPC = 0.75 and treat the price level as fixed.

\begin{enumerate}
    \item[(a)] Compute the expenditure multiplier. (1 mark)
    \item[(b)] Calculate the total change in real GDP. (2 marks)
    \item[(c)] Compute the new real GDP. (1 mark)
    \item[(d)] Explain AD--AS short-run effects on GDP, unemployment, and price level. (3 marks)
\end{enumerate}


\vspace{5em}

\textbf{Problem 2 (8 marks) — Simple Deposit Multiplier}

A new deposit of \$5,000 is placed into a Canadian commercial bank.  
The desired reserve ratio is:
\[
r_d = 0.10.
\]

\begin{enumerate}
    \item[(a)] Compute the simple deposit multiplier. (1 mark)
    \item[(b)] Compute the \textit{maximum} possible increase in deposits. (2 marks)
    \item[(c)] From the initial deposit, how much must be held as reserves? How much can be loaned? (3 marks)
    \item[(d)] Give one reason the actual deposit expansion may be smaller. (2 marks)
\end{enumerate}

\hrule
\vspace{1em}

\begin{center}
\textit{End of Quiz}
\end{center}

\end{document}
